\documentclass{article}

\usepackage[utf8]{inputenc}
\usepackage[T1]{fontenc}
\usepackage{amsmath,amssymb,amsthm}
\usepackage{mathtools}
\usepackage{hyperref}
\usepackage{booktabs}

\makeatletter
\newcommand{\citep}[2][]{%
  \begingroup
  \def\@tmpkey{#2}%
  \ifx\@tmpkey\@empty\else
    (\ifx\@empty#1\@empty\else #1, \fi
     \@nameuse{bibkey@\@tmpkey})%
  \fi
  \endgroup
}
\expandafter\def\csname bibkey@lasser2026horizon\endcsname{Horizon~Aware}
\expandafter\def\csname bibkey@lasser2026exo\endcsname{Exogenous~Verification}
\expandafter\def\csname bibkey@lasser2026phase\endcsname{Phase~Redundancy}
\makeatother

\newtheorem{definition}{Definition}
\newtheorem{lemma}{Lemma}
\newtheorem{theorem}{Theorem}
\newtheorem{remark}{Remark}
\newtheorem{assumption}{Assumption}

\newcommand{\rext}{r_{\mathrm{ext}}}
\newcommand{\meff}{m_{\mathrm{eff}}}
\newcommand{\mindep}{m_{\mathrm{eff}}^{\mathrm{indep}}}
\newcommand{\Saudit}{S_*}

\title{Lemma 5a (Substrate floor): Full Proof Draft v2}
\author{Paper 10 working draft for codex re-review}
\date{}

\begin{document}

\maketitle

\section{Goal and changes from v1}

V1 incorrectly claimed Step~4's rate additivity followed from joint
failure-correlation independence ($I_6'$).  Codex Round~A flagged
this as a significant problem: $I_6'$'s definition only says
adversarial shock mechanisms are single-substrate-supported, which
does \emph{not} imply cooperative-production channels are
statistically independent or non-overlapping.  V2 fixes this and
incorporates Round~A's other refinements.

\textbf{Changes from v1:}
\begin{itemize}
\item \textbf{Q1 fix.}  Use heterogeneous per-pair rates
  $\rho_{ij}$ with $\rho_0 := \min_{i<j \in \Saudit} \rho_{ij}$
  over an \emph{audited subset} $\Saudit$ of $m^*$ substrates,
  not every pair in the deployment population.
\item \textbf{Q2 fix (the big one).}  Replaced the incorrect
  ``joint independence implies additivity'' argument with a new
  explicit structural assumption: \textbf{(SU) Pairwise cooperative
  channel superposition.}  Exact-two-substrate cooperative event
  classes are disjointly attributed and not rivalrous across pairs
  except through explicitly-modeled higher-order channels.
\item \textbf{Q3 fix.}  Pairwise channels defined as
  \textbf{exact two-substrate support}, not ``at least one
  participant from each.''  This eliminates double-counting of
  higher-order cooperatives in the pairwise sum.
\item \textbf{Q5 fix.}  CN-with-superposition surfaced as theorem-
  level condition (C10) of Theorem~\ref{thm:deployment_safety},
  analogous to (C6)--(C9) added in earlier integrations.
\item \textbf{Caveat fix.}  Removed v1's incorrect claim that
  failure-correlation independence excludes correlated-channel
  sub-additivity.  Failure-correlation independence excludes
  shared shock mechanisms; (SU) is a separate structural condition
  that excludes cooperative-production overlap.
\end{itemize}

\section{Source-paper grounding}

\citep[Horizon~Aware]{lasser2026horizon}'s anti-monopolar property
establishes that under a discounted objective with $\rext > 0$,
domination is anti-maximizing.  Paper~3 grounds $\rext > 0$ in
substrate physics (cross-substrate cooperatives produced by
distinct-substrate participants whose joint output cannot be
replicated by any single substrate) and shows pairwise-novel
cooperative output between any pair of distinct substrates.
However, Paper~3 does not give a quantitative floor on $\rext$ as a
function of substrate count.

Lemma~5a's contribution is to derive an explicit lower bound on
$\rext$ given (a)~the substrate count $\mindep$, (b)~per-pair
calibrated cooperative novelty rates, and (c)~a pairwise-channel
superposition structural condition.

\section{Pairwise cooperative channels (exact-two-substrate support)}

\begin{definition}[Pairwise cooperative channel]
\label{def:pairwise_channel}
For an audited subset $\Saudit \subseteq \mathcal{S}$ of $m^*$
distinct substrates from the deployment's substrate population
$\mathcal{S}$, the \emph{pairwise cooperative channel} for the pair
$(s_i, s_j)$ with $i < j$ is the event class:
\[
  \mathcal{C}^{(s_i, s_j)} \;:=\; \{c \in P :
  \mathrm{participants}(c) \subseteq s_i \cup s_j,\
  \mathrm{participants}(c) \cap s_i \neq \emptyset,\
  \mathrm{participants}(c) \cap s_j \neq \emptyset\},
\]
i.e., cooperatives produced \emph{exactly} by participants from
$s_i$ and $s_j$ (no participants from other substrates).
\end{definition}

\textbf{Note (Q3 fix).}  V1's definition admitted ``cooperatives
whose participants \emph{include at least one} from $s_i$ and
$s_j$,'' which over-counts higher-order cooperatives across multiple
pairs.  V2's exact-two-substrate-support definition prevents this
double-counting; higher-order cooperatives belong to separate
event classes and contribute to $\rext$ additionally
(Remark~\ref{rem:higher_order}).

\section{Coercivity and superposition assumptions}

\begin{assumption}[CN: Per-pair cooperative-novelty rate floor, heterogeneous form]
\label{ass:CN}
For an audited subset $\Saudit \subseteq \mathcal{S}$ with $|\Saudit|
= m^*$ and every pair $(s_i, s_j)$ with $i < j$ in $\Saudit$:
\[
  \rext^{(s_i, s_j)} \;\geq\; \rho_{ij} \;>\; 0,
\]
where $\rext^{(s_i, s_j)}$ is the per-pair contribution to $\rext$
from the pairwise cooperative channel $\mathcal{C}^{(s_i, s_j)}$,
and $\rho_{ij}$ are per-pair rate floors.  Define $\rho_0 :=
\min_{i<j \in \Saudit} \rho_{ij} > 0$.

The constants $\rho_{ij}, \rho_0$ depend on the deployment class
$\mathcal{D}$ but not on $|P|$.
\end{assumption}

\begin{assumption}[SU: Pairwise cooperative channel superposition]
\label{ass:SU}
For the audited subset $\Saudit$ and the pairwise cooperative
channels $\{\mathcal{C}^{(s_i, s_j)}\}_{i<j}$ of
Definition~\ref{def:pairwise_channel}:
\begin{enumerate}
\item \textbf{Disjoint attribution:} the channels are pairwise
  disjoint as event classes: $\mathcal{C}^{(s_i, s_j)} \cap
  \mathcal{C}^{(s_k, s_l)} = \emptyset$ for $\{i,j\} \neq \{k,l\}$.
\item \textbf{Non-rivalrous production:} the rate of cooperative
  production in channel $\mathcal{C}^{(s_i, s_j)}$ is not reduced
  by simultaneous production in other pairwise or higher-order
  channels (no rivalrous capability allocation across pairs).
\item \textbf{Joint-deployment intensities:} the per-pair rates
  $\rho_{ij}$ are measured in the joint deployment with all
  audited-subset substrates active, not under a counterfactual
  isolation.
\end{enumerate}
\end{assumption}

\textbf{Operational interpretation of SU.}  SU is the structural
condition that allows additive aggregation of per-pair rates.
Without SU, two substrate pairs could share cooperative
``novelty bottlenecks'' (e.g., a finite supply of joint-attention
slots between $s_i$ agents and other substrates), in which case
producing cooperatives in channel $\mathcal{C}^{(s_i, s_j)}$ would
reduce production in channel $\mathcal{C}^{(s_i, s_k)}$.  SU
asserts that such rivalrous coupling is absent in the audited
subset.

SU is a strictly stronger assumption than $I_6'$'s joint failure-
correlation independence: $I_6'$ excludes shared adversarial shock
mechanisms; SU excludes shared cooperative-production bottlenecks.
The two conditions are structurally independent.

\section{Statement of Lemma 5a}

\begin{lemma}[Substrate floor, v2]
\label{lem:5a}
Under invariant $I_6'$ ($\mindep \geq m^* \geq 3$), condition
(C10) of Theorem~\ref{thm:deployment_safety} (per-pair cooperative-
novelty rate floor with superposition; comprising
Assumptions~\ref{ass:CN} and~\ref{ass:SU} on an audited substrate
subset $\Saudit$):
\[
  \rext \;\geq\; r_*(m^*) \;:=\; \rho_0 \cdot \binom{m^*}{2} \;>\; 0.
\]
For $m^* = 3$, $r_*(3) = 3\rho_0 > 0$.

The floor $r_*(m^*)$ is intensive in $|P|$ over the deployment
class $\mathcal{D}$: $\rho_0$ depends on $\mathcal{D}$'s operational
parameters but not on $|P|$.
\end{lemma}

\section{Proof}

\begin{proof}
We work over a deployment class $\mathcal{D}$ satisfying $I_6'$ and
(C10) (Assumptions CN, SU on audited subset $\Saudit$).

\emph{Step 1 (substrate audit).}  By $I_6'$, the deployment has
$\mindep \geq m^*$ substrates with joint failure-correlation
independence.  By (C10), an audited subset $\Saudit \subseteq
\mathcal{S}$ with $|\Saudit| = m^*$ has been verified to satisfy
both CN (per-pair rate floor) and SU (cooperative channel
superposition).  We work over $\Saudit$.

\emph{Step 2 (combinatorial pairwise channels).}  The number of
unordered pairs in $\Saudit$ is $\binom{m^*}{2} = m^*(m^*-1)/2$.
Each pair $(s_i, s_j)$ with $i < j$ has its own pairwise
cooperative channel $\mathcal{C}^{(s_i, s_j)}$ (exact-two-substrate
support, Definition~\ref{def:pairwise_channel}).

\emph{Step 3 (per-pair rate from CN).}  By Assumption~CN,
$\rext^{(s_i, s_j)} \geq \rho_{ij} \geq \rho_0$ for every pair in
$\Saudit$.

\emph{Step 4 (additivity from SU).}  By Assumption~SU(1) (disjoint
attribution), the pairwise channels do not double-count
cooperatives.  By Assumption~SU(2) (non-rivalrous production), the
realized rate from each channel does not depend on production in
other channels.  By Assumption~SU(3) (joint-deployment
intensities), the per-pair rates $\rho_{ij}$ measure production
in the actual deployment, not in counterfactual isolation.

Therefore, the pairwise contributions sum without overcounting or
suppression:
\[
  \rext^{\mathrm{pairwise}} \;:=\; \sum_{i<j \in \Saudit}
  \rext^{(s_i, s_j)} \;\geq\; \binom{m^*}{2} \cdot \rho_0.
\]
The total $\rext$ includes pairwise plus any higher-order
cooperative contributions; pairwise is therefore a lower bound:
\[
  \rext \;\geq\; \rext^{\mathrm{pairwise}} \;\geq\; \binom{m^*}{2}
  \cdot \rho_0 \;=\; r_*(m^*).
\]

\emph{Step 5 (intensivity).}  $\rho_0 = \min_{i<j \in \Saudit}
\rho_{ij}$ is a deployment-class constant by Assumption~CN
(per-pair calibrated rates), independent of $|P|$.  $m^*$ is the
$I_6'$ threshold, also a deployment-class constant.  Therefore
$r_*(m^*) = \rho_0 \binom{m^*}{2}$ is independent of $|P|$ over
$\mathcal{D}$.  $\Box$
\end{proof}

\section{Remarks on scope}

\begin{remark}[On the audited subset $\Saudit$]
\label{rem:saudit}
v1 used ``every pair in the deployment population.''  v2 uses an
audited subset $\Saudit$ of exactly $m^*$ substrates.  Reason:
deployments may have weak substrate pairs (e.g., recently-added
substrates that haven't yet reached calibrated cooperative
novelty rates).  Including all pairs would force $\rho_0$ to drop
to the weakest pair, weakening the bound.

The $\Saudit$ formulation lets the deployment audit the strongest
$m^*$ substrates and bound $\rext$ from below on that subset.
$\rext$ for the full deployment is at least the audited-subset
contribution; weak substrate pairs contribute non-negatively but
are not credited.

Operationally, the audit (Audit~6 of
\S\ref{sec:tooling_substrate_audits}) calibrates $\rho_{ij}$ for
each candidate pair, picks the $m^*$ substrates whose pairwise
rates collectively maximize $\binom{m^*}{2} \cdot \min \rho_{ij}$.
\end{remark}

\begin{remark}[On higher-order cooperatives]
\label{rem:higher_order}
The lemma's pairwise-only bound is conservative.  Higher-order
cooperatives (those with exact $k$-substrate support for $k \geq
3$) contribute additional rate to $\rext$:
\[
  \rext \;\geq\; \rext^{\mathrm{pairwise}} +
  \sum_{k=3}^{m^*} \rext^{(k\text{-way})}.
\]
The $k$-way contributions can be calibrated analogously
(per-$k$-tuple rate floors $\rho_0^{(k)}$ and superposition
conditions), giving a tighter floor:
\[
  \rext \;\geq\; \sum_{k=2}^{m^*} \rho_0^{(k)} \binom{m^*}{k}.
\]
For deployment-claim purposes, the pairwise floor suffices: it
makes Layer~1's safe region non-trivial and is intensive in $|P|$.
Operators may use higher-order calibration for tighter bounds.
\end{remark}

\begin{remark}[On the SU assumption's strength relative to $I_6'$]
\label{rem:SU_vs_I6}
v1 incorrectly stated that joint failure-correlation independence
($I_6'$, Definition~\ref{def:joint_indep}) excludes correlated-
channel sub-additivity.  Codex Round~A correctly flagged this:
$I_6'$ excludes \emph{shared adversarial shock mechanisms}
(``no single shock spans multiple substrates''), not \emph{shared
cooperative-production bottlenecks} (``no rivalrous resource
allocation across pairs'').  These are structurally independent
conditions.

(SU) is a new structural condition added by Paper~10 specifically
for Lemma~5a.  Failure-correlation independence and superposition
are both required; both are new theorem-level conditions under
(C2) and (C10) respectively.
\end{remark}

\section{Connection to Theorem 1 Layer 1}

Lemma~5a's role: provide the substrate-floor input that makes
Lemma~5c's safe-region inequality non-trivial.

Lemma~5c asserts $\Delta r_K < \rext + (1-\gamma)(\Ddiv -
\Delta_0)$ implies diversity dominates.  For non-vacuous safe
region, $\rext$ must be bounded below by a verifiable constant.
Lemma~5a provides $\rext \geq \rho_0 \binom{m^*}{2} > 0$ under
(C10), so the safe region is non-trivial.

Lemma~5c's appendix proof Step~1 invokes Lemma~5a:
\[
  \rext \geq r_*(m^*) > 0 \quad \text{(by Lemma~\ref{lem:5a})},
\]
which guarantees the safe-region inequality has positive right-hand
side under any positive $\Delta_0$ and $\gamma < 1$.

\section{What this proof does and does not establish}

\textbf{Establishes:}
\begin{itemize}
\item A formal coercivity assumption (CN: per-pair cooperative-
  novelty rate floor with audited subset $\Saudit$ and
  heterogeneous rates $\rho_{ij}$).
\item A formal superposition assumption (SU: disjoint attribution,
  non-rivalrous production, joint-deployment intensities) replacing
  v1's incorrect ``joint independence implies additivity'' argument.
\item Pairwise-channel definition with exact-two-substrate support
  (no double-counting).
\item Combinatorial-additivity proof of the $\binom{m^*}{2}$
  scaling under SU.
\item Closed-form floor $r_*(m^*) = \rho_0 \binom{m^*}{2}$
  intensive in $|P|$.
\item Three remarks: $\Saudit$ scope, higher-order cooperatives,
  SU vs.\ $I_6'$ independence.
\end{itemize}

\textbf{Does not establish (deferred):}
\begin{itemize}
\item Tightness of $r_*(m^*)$.  Pairwise-only floor is conservative;
  higher-order cooperatives contribute additionally
  (Remark~\ref{rem:higher_order}).
\item Existence of deployment configurations satisfying CN, SU.
  Conditions on the deployment's substrate structure;
  verification is empirical (Audit~6 of
  \S\ref{sec:tooling_substrate_audits}, to be added).
\item Numerical values of $\rho_{ij}, \rho_0$.  Calibrated
  deployment-by-deployment via per-pair cooperative-novelty
  rate-estimation procedures.
\end{itemize}

\section{Specific verification questions for v2 codex review}

\begin{enumerate}
\item Is the SU formulation (disjoint attribution, non-rivalrous
  production, joint-deployment intensities) the right structural
  condition for additivity?  Are the three sub-conditions all
  necessary, or can some be derived from others?

\item Is the audited-subset $\Saudit$ formulation correctly
  scoped?  In particular, is it an audit-time construct (operators
  choose $\Saudit$ for the calibration) or a deployment-state
  construct (the $m^*$-best substrates may change over deployment)?

\item Is the exact-two-substrate-support pairwise-channel
  definition correctly preventing double-counting?  In particular,
  if a deployment has cooperatives with two-substrate support but
  some optional third-substrate participants (sometimes producing
  3-way cooperatives, sometimes 2-way), how should these be
  classified?

\item Is Step~4's additivity argument rigorous under v2's SU?
  The v1 mistake was conflating $I_6'$ with cooperative-channel
  independence.  Does v2's SU now correctly cover what's needed,
  or are there further structural conditions hiding?

\item Is the (C10) framing for Theorem 1 correct?  Specifically:
  CN + SU together as a single condition, or as two separate
  conditions (C10) and (C11)?
\end{enumerate}

\end{document}
